%% -------------------------------------------------------
%% Anhang A: API-Spezifikation
%% -------------------------------------------------------
\chapter{API-Endpunkte im Detail}
\label{app:api}

Dieser Anhang listet alle implementierten und geplanten REST-Endpunkte der
DoggoGO-Backend-API mit Request/Response-Schemata auf. Die vollständige und interaktive
Dokumentation ist über die Swagger-UI unter \texttt{/swagger} erreichbar.

\section*{Authentifizierung}

\textbf{POST /api/owners/register}

\begin{itemize}
  \item \textbf{Beschreibung:} Registriert eine neue Nutzerin / einen neuen Nutzer.
  \item \textbf{Request-Body (JSON):}
\end{itemize}
\begin{lstlisting}[language={}]
{
  "name": "Max Muster",
  "email": "max@example.com",
  "password": "SecurePass123!",
  "phoneNumber": "+43 123 456789"
}
\end{lstlisting}
\begin{itemize}
  \item \textbf{Response 201:} Owner-Objekt (ohne Passwort-Felder)
  \item \textbf{Response 409:} E-Mail bereits vergeben
\end{itemize}

\vspace{1em}
\textbf{POST /api/owners/login}
\begin{itemize}
  \item \textbf{Beschreibung:} Meldet eine Nutzerin / einen Nutzer an und gibt ein JWT zurück.
  \item \textbf{Request-Body (JSON):}
\end{itemize}
\begin{lstlisting}[language={}]
{
  "email": "max@example.com",
  "password": "SecurePass123!"
}
\end{lstlisting}
\begin{itemize}
  \item \textbf{Response 200:} \texttt{\{"token": "eyJ...\}"}
  \item \textbf{Response 401:} Ungültige Anmeldedaten
\end{itemize}

\section*{Hunde (Dogs)}

\textbf{GET /api/dogs} \quad \textit{(Authentifizierung erforderlich)}\\
Gibt alle Hundeprofile der aktuellen Nutzerin / des aktuellen Nutzers zurück.

\textbf{POST /api/dogs} \quad \textit{(Authentifizierung erforderlich)}\\
Legt ein neues Hundeprofil an.
\begin{lstlisting}[language={}]
{
  "name": "Bello",
  "age": 3,
  "breedId": "a1b2c3d4-..."
}
\end{lstlisting}

\textbf{PUT /api/dogs/\{id\}} \quad \textit{(Authentifizierung erforderlich)}\\
Aktualisiert ein bestehendes Hundeprofil.

\textbf{DELETE /api/dogs/\{id\}} \quad \textit{(Authentifizierung erforderlich)}\\
Löscht ein Hundeprofil.

\section*{Hunderassen (Breeds)}

\textbf{GET /api/breeds}\\
Gibt alle Hunderassen inklusive \texttt{LegalRestrictions} zurück.

\textbf{GET /api/breeds/\{id\}}\\
Gibt eine einzelne Rasse mit vollständigen Rechtsvorschriften zurück.
